%Luchando los dias 
%============================================
%  Explicar Los daots que se usaran para este promblema 
%=============================================

\begin{frame}{Generación de Datos Simulados}

    \begin{columns}[c]

        %========================================
        % COLUMNA IZQUIERDA: Tabla de Sensores
        %========================================
        \column{0.50\textwidth}
        \begin{block}{Ubicación de Sensores ($S_i$)}
            \centering
            \scriptsize 
            \begin{tabular}{@{}lccc@{}}
                \toprule
                \textbf{Sensor} & \textbf{$x_i$} & \textbf{$y_i$} & \textbf{$z_i$} \\ \midrule
                S1 & -4.0 & -3.0 & 0.2 \\
                S2 & -2.5 & 2.8 & -0.1 \\
                S3 & 0.0 & 4.0 & 0.1 \\
                S4 & 3.5 & 3.0 & -0.2 \\
                S5 & 4.5 & -1.5 & 0.0 \\
                S6 & 2.0 & -4.0 & 0.3 \\
                S7 & -1.0 & -4.5 & -0.1 \\
                S8 & -4.5 & 1.0 & 0.2 \\
                S9 & 0.5 & 0.5 & 0.0 \\
                S10 & 2.8 & 0.8 & -0.3 \\
                S11 & -1.8 & -1.2 & 0.1 \\
                S12 & 1.0 & 3.2 & 0.2 \\ \bottomrule
            \end{tabular}
        \end{block}

        %========================================
        % COLUMNA DERECHA: Resumen Ejecutivo
        %========================================
        \column{0.45\textwidth}

        \begin{block}{Configuración de Red}
            \small
            \begin{itemize}
                \item \textbf{Muestra:} $n = 12$ sensores.
                \item \textbf{Espacio:} Distribución tridimensional.
            \end{itemize}
        \end{block}

        \vspace{0.4cm}

        \begin{block}{Estrategia de Simulación}
            \scriptsize
            \begin{itemize}
                \item Cobertura multidireccional.
                \item Estabilidad del problema inverso.
                \item Generación de amplitudes ideales.
            \end{itemize}
        \end{block}

        \vspace{0.6cm}
        
        \begin{center}
            \tiny \textcolor{gray}{Red optimizada para localización robusta de la fuente.}
        \end{center}

    \end{columns}

\end{frame}